% s13 The main theorem for King's printed state
\section{The main theorem for King's printed state}

We state the main theorem in full and prove it by assembling Sections~4--12. The object is King's printed state $\ZKinf(f,0)=\lim_J\lim_k\avK{e^{F(f)}}$ on King's RG cube (\cref{def:s02:Kcube}), block factor $5$, at a fixed photon mass $\lmu>0$ and all positive couplings $e,\lambda,m$ (\cref{lem:s14:CW}). The curvature sector (Leg~C) is proved on King's convention $\NDK$. The Wick sector (Leg~W) is proved in its tube form, Form~(B), on the comparison convention $\NDa$ (\cref{def:s02:acube}), on which \cref{thm:s11:W,lem:s11:TWtube,thm:s08:OS2WT} are written and whose functionals $Z^a_{k,J}$ are the ones of \cref{hyp:s03:W1}; it is carried to King's convention by \cref{thm:s13:TWK,prop:s13:WK}: the two tube statements that the interlacing of \cref{lem:s04:G} does not transfer, the bound \eqref{eq:s11:W} and the fixed-cube limits on the tube for the K-family functionals $Z^{a,K}_{k,J}$, are proved there directly. The torus inputs are \cref{hyp:s03:T1,hyp:s03:T2}, and the coupling condition is the sign clause \Tthree\ (\cref{hyp:s03:T3}), proved on the explicit regions of \cref{prop:s14:T3}; $\cK=\{(\mathrm if,g):g\ge0\}$ is the cone. \Cref{thm:s01:main} is the theorem below.

\begin{theorem}[the axioms for King's printed state at fixed photon mass]\label{thm:s13:main}
Let $e>0$, $\lambda,m,\lmu>0$, $\eps_k=5^{-k}$, $\CJ=[-MJ,MJ)^3$ with King's block integer $M\ge1$, and let $\NDK_{k,J}$ be King's cube state.

\smallskip\noindent\textbf{A. The object.} For every finite family $f=(f_1,\dots,f_N)$ of real $2$-forms in $C^\infty_c(\R^3;\Lambda^2)$ the limit
\[
\ZKinf\Big(\sum_jz_jf_j,0\Big)=\lim_{J\to\infty}\lim_{k\to\infty}\avK{e^{\sum_jz_jF_{\Lk}(f_j)}}
\]
exists locally uniformly in $z\in\C^N$, is entire, depends on the family only through the total source, is real and $\ge1$ at real $z$, lies in $[0,1]$ at purely imaginary $z$, and satisfies $|\ZKinf(\sum_jz_jf_j,0)|\le e^{\norm{g(z)}_2^2}$, $g(z)=\sum_j(\Re z_j)f_j$, with no dependence on $\lmu$, $e$, $\lambda$, $m$ or $M$.

\smallskip\noindent\textbf{B. Leg~C (the curvature field).} The Schwinger functions $S^K_n(f_1,\dots,f_n)=\partial_{z_1}\cdots\partial_{z_n}\ZKinf(\sum_jz_jf_j,0)|_{z=0}$ satisfy:
\begin{enumerate}[label=\textup{(\roman*)},leftmargin=*]
\item \emph{\OSz.} $|\ZKinf(\sum_jz_jf_j,0)|\le e^{\norm{g(z)}_2^2}$ for all $z\in\C^N$, and also $\le\exp(\frac12\lmu^{-2}\norm{\sum_j(\Re z_j)\dstar f_j}_2^2)$; the same two bounds, with the lattice norms $\norm{\cdot}_{2,k}$ and $\dstar_k$, hold at every $(k,J)$ for the lattice functionals $\avK{e^{\sum_jz_jF_{\Lk}(f_j)}}$. For real test functions
\[
|S^K_n(f_1,\dots,f_n)|\le(\sqrt2\,e)^n(n!)^{1/2}\prod_{i=1}^n\norm{f_i}_2\le(e\pi/2)^n(n!)^{1/2}\prod_{i=1}^n\fournorm{f_i};
\]
each $S^K_n$ extends to a continuous symmetric multilinear form on $\cS(\R^3;\Lambda^2_\C)^n$, a tempered distribution, the bound carrying the factor $2^n$ for complex test functions; so \Ezz\ holds with $s=4$, $\sigma_n=(e\pi)^n(n!)^{1/2}$ and every $\beta>\frac12$.
\item \emph{\OSo.} The $S^K_n$ are invariant under all translations of $\R^3$ and under the cubic group $\Oh$ about every point, $F$ transforming as a $2$-form. Moreover $\ZKinf$ is invariant under the Euclidean group $\Euc$ and the $S^K_n$ are Euclidean covariant.
\item \emph{\OSt.} Reflection positivity about every coordinate plane $\{x_\nu=c\}$, $\nu\in\{1,2,3\}$, $c\in\R$: for finite families $\zeta_1,\dots,\zeta_N$ of real $2$-form sources supported in $\{x_\nu>c\}$ with arbitrary complex coefficients, $\sum_{i,j}\bar\gamma_i\gamma_j\ZKinf(\theta\zeta_i^*\sqcup\zeta_j,0)\ge0$ for all $\gamma\in\C^N$ (the matrix form for exponentials), and $S(\theta(P)P)\ge0$ for every polynomial $P$ in the smeared curvature with sources supported in $\{x_\nu>c\}$ (the polynomial form). The same holds about every plane $\{x\cdot n=c\}$, $n$ a unit vector.
\item \emph{\OSth.} $S^K_n$ is real for real test functions, symmetric, and zero for odd $n$.
\item \emph{The measure.} For every even $r\ge14$ there is exactly one Borel probability measure $\mu^{K,H}$ on $\cH^{-r}(\R^3;\Lambda^2)$ whose moments on $C^\infty_c$ are the $S^K_n$; its characteristic functional is $\ZKinf(\mathrm i\,\cdot,0)$, extended by continuity to $\cS$; its image $\mu^K_S$ on $H^{-r}_{\rm loc}(\R^3;\Lambda^2)$ is the unique Borel probability measure there with this characteristic functional, and the unique one with these moments. For every real $\chi\in C^\infty_c(\R^3)$ and every box $Q$ containing $\supp\chi$ in its interior, the laws of the localised lattice curvature fields $\chi\Phi^K_{k,J}$ under $\NDK_{k,J}$ on $H^{-r}_{\rm per}(Q;\Lambda^2)$ converge weakly, along the full sequence $k\to\infty$ at every fixed $J\ge J_0(\chi)$ and then as $J\to\infty$, to the law of $\chi\Phi$ under $\mu^K_S$. Nothing is asserted about the laws of the unlocalised lattice fields at fixed $J$.
\item \emph{Weak \OSf.} $\mu^K_S$ is $\Cub$-invariant and is an integral $\int\mu_\xi\,d\nu(\xi)$ of translation-ergodic, reflection-positive, $\theta_j$-invariant components which cluster in the mean and under diagonal separation, have a unique joint vacuum in each coordinate direction, and satisfy the factorial growth bound $|S^\xi_N(u_1,\dots,u_N)|\le\sigma^\xi_N\prod_i\norm{u_i}_{r+2}$, $\sigma^\xi_N\le\alpha_\beta(\xi)(N!)^\beta$ for every $\beta>1$, with component constants (no tube bound); there is an $\Oh$-invariant decomposition into mixtures of $k\in\{1,2,3,6\}$ ergodic measures; and $\nu$-almost every component satisfies \OSf\ in the weak form (ergodicity under $\R^3$). Moreover $\nu$-almost every component is $\Euc$-invariant, reflection positive about every plane, and satisfies \OSf\ in each of its forms for every time direction, with $\ker H_n=\C\Omega_n$ in every direction $n$; its moments are real and symmetric. The vanishing of a component's odd moments is not claimed.
\end{enumerate}

\smallskip\noindent\textbf{C. Leg~W, Form~(B), on the comparison convention.} Let the couplings satisfy the sign clause \Tthree\ (\cref{hyp:s03:T3}; in particular, let them lie in one of the explicit regions of \cref{prop:s14:T3}), and fix an admissible scheme $a$ (\cref{def:s03:adm}; $z=a^\sharp$ the reference scheme). Let $\sigma_+=s_1$, $H=s_1/16$, $\kappa=C_U/s_1$, $K_-$ be the constants of \eqref{eq:s11:constants} ($K_-$ built from $B_{\rm cell}$; for $a\ne z$ with $C_U$ replaced by $C_U+s_1\sup_k|z_k-a_k|$), $\sigma_0=H/4$ and $K'=\kappa\pi^2+K_-$. Then for every finite family $\cF=(f_1,\dots,f_n;h_1,\dots,h_p)$ of real $2$-forms and real scalars the tube object
\[
\Ginf(\zeta)=\lim_{J\to\infty}\lim_{k\to\infty}Z^a_{k,J}\Big(\sum_iz_if_i,\sum_jw_jh_j\Big),\qquad\zeta=(z,w)\in\Dtube=\{\norm{k_+(\zeta)}_\infty<\sigma_+/2,\ \norm{k_-(\zeta)}_\infty<H/2\},
\]
exists locally uniformly on the tube and is holomorphic there; on the Wick tube of \cref{def:s08:enlarged} the limit $Z^a_{\infty,T}$ of \cref{prop:s08:T} exists and equals $\Ginf$ where both are defined. The joint Schwinger functions $S_{n,p}(f_1,\dots,f_n;h_1,\dots,h_p)=\partial^{(n)}_z\partial^{(p)}_w\Ginf(0)$ of the pair $(F,W^a)$ satisfy:
\begin{enumerate}[label=\textup{(\roman*)},leftmargin=*]
\item \emph{\OSz.} $|\Ginf(\zeta)|\le\exp(\norm{g(\zeta)}_2^2+\kappa\norm{k_+(\zeta)}_{L^1}+K_-\wnorm{k_-(\zeta)})$ on $\Dtube$, and the same bound with lattice norms holds for $Z^a_{k,J}$ uniformly in $k\ge0$ and $J\ge J_1(R)$ when the supports lie in $Q_R$; for real test functions
\[
|S_{n,p}(f_1,\dots,f_n;h_1,\dots,h_p)|\le e^{K'\sigma_0}(\sqrt2\,e)^n(n!)^{1/2}(e/\sigma_0)^p\,p!\prod_{i=1}^n\norm{f_i}_2\prod_{j=1}^p\wnorm{h_j},
\]
with the factor $2^{n+p}$ for complex test functions; the $S_{n,p}$ are tempered distributions and \Ezz\ holds with $s=4$ and every $\beta>1$. The constants $\sigma_+,H,\kappa,K_-$ depend on $(e,\lambda,m,\lmu,M,a)$, on the fixed auxiliary choices $M_1$, $\chi_0$ and the amplitude $s_1\in(0,\sigma_1)$ of \Itwo, and on nothing else; they are not uniform in $\lmu$.
\item \emph{\OSo.} $\Ginf(\tau_bR\cdot\zeta)=\Ginf(\zeta)$ on $\Dtube$ for every $b\in\R^3$ and $R\in\Oh$, $F$ transforming as a $2$-form and $W^a$ as a scalar (the transformed family has the same tube); hence the $S_{n,p}$ are invariant under translations and under $\Oh$ about every point. No Euclidean clause is asserted.
\item \emph{\OSt.} For finite families $\zeta_1,\dots,\zeta_N$ of real sources supported in $\{x_\nu>c\}$ with complex coefficients of amplitude $\amp(\zeta_i)<\sigma_W\defeq\sigma_+/8$ (in particular all purely imaginary Wick coefficients with arbitrary complex curvature coefficients), $\sum_{i,j}\bar\gamma_i\gamma_jZ^a_{\infty,T}(\theta\zeta_i^*\sqcup\zeta_j)\ge0$ about every coordinate plane; and the polynomial form $S(\theta(P)P)\ge0$ for polynomials in the smeared $F$ and $W^a$ with sources in $\{x_\nu>c\}$.
\item \emph{\OSth.} The $S_{n,p}$ are symmetric multilinear in the $f$'s and in the $h$'s, real for real test functions, and zero for odd $n$.
\item[\textup{(v),(vi)}] \emph{The measure and weak \OSf.} Clauses (v) and (vi) of Leg~C hold for the joint field on $\cH^{-r}(\R^3;\Lambda^2\oplus\R)$ and $H^{-r}_{\rm loc}(\R^3;\Lambda^2\oplus\R)$, every even $r\ge14$, with moments $S_{n,p}$, characteristic functional $(u,v)\mapsto\Ginf(\mathrm iu,\mathrm iv)$ on real test functions (purely imaginary coefficients have $k(\zeta)=0$ and lie in the tube), the component growth bound with the constants of (i), and the local convergence of the laws of the scheme-shifted lattice fields $\chi(F_{\Lk},W^a_{\Lk})$ along the full sequence $k\to\infty$ at every fixed $J$ with $MJ>R(\chi)+\ell_c$, then $J\to\infty$. No Euclidean clause is asserted. Two admissible schemes $a,a'$ give objects differing by the shift $W\mapsto W-\kappa_{a'a}\int(\cdot)$ of the Wick component, $\kappa_{a'a}=\lim_k(a'_k-a_k)$; joint cumulants of order $\ge2$ are scheme-free.
\end{enumerate}
No bound at positive scalar sources beyond the cone $\sup q\le s_1$ of \cref{cor:s11:W2prime} is assumed. The arbitrary-amplitude form (Form~(A)), which carries the only Euclidean clause of the Wick sector, is the optional extension of \cref{thm:s23:FormA}.

\smallskip\noindent\textbf{D. Leg~W on King's convention.} Under the conditions of~C, clauses (i)--(vi) of~C hold, with the same constants, for the tube object $G^K_\infty=\lim_J\lim_kZ^{a,K}_{k,J}$ of King's printed state on $\NDK$ (\cref{thm:s13:TWK,prop:s13:WK}).
\end{theorem}

\begin{proof}
\emph{A.} \Cref{thm:s05:CK}(b): by \cref{hyp:s03:EC} the two limits exist locally uniformly on $\C^N$ and are entire, $|\ZKJ|,|\ZKinf|\le e^{\norm{g(z)}_2^2}$, $J\mapsto\ZKJ(\mathrm it)$ is nondecreasing with values in $[0,1]$, and $\ZKinf(x)\ge1$ at real $x$; realness at real $z$ is the symmetry $(A,\vartheta)\mapsto(-A,-\vartheta)$ of every $\NDK_{k,J}$, closed under limits. Each lattice functional depends on the family only through its total source, hence so do the limits. The bound is absolute (\cref{rem:s05:absolute}).

\emph{B.} (i) is \cref{thm:s05:CK}: the finite-volume bound (a) with $\min\{\norm{g}^2_{2,k},\frac12\lmu^{-2}\norm{\dstar_kg}^2_{2,k}\}$ at every $(k,J)$; the two limit bounds (b), the second being (b)'s second bound, carried from the lattice by $\norm{\dstar_kg}_{2,k}\to\norm{\dstar g}_2$, \eqref{eq:s02:riemann}; the Schwinger bounds, the extension to $\cS$, the factor $2^n$ and \Ezz\ (c). (ii): translations are \cref{cor:s06:TK} (the lattice sandwich \cref{lem:s06:TK}, the two limits, the identity theorem, density); the cubic group about every point is \cref{thm:s06:S1K}; the Euclidean clause is \cref{thm:s06:O3}, that is the rotation theorem \cref{hyp:s03:RC} with \cref{cor:s04:G2} and \cref{cor:s06:TK}. (iii): the matrix form about every coordinate plane and the polynomial form are \cref{lem:s08:OS2K} with \cref{prop:s08:poly}; a plane $\{x\cdot n=c\}$ is the image of a coordinate plane under a Euclidean motion $\gamma$, and by (ii) $\ZKinf(\theta_\gamma\zeta_i^*\sqcup\zeta_j,0)=\ZKinf(\gamma^{-1}\cdot(\theta_\gamma\zeta_i^*\sqcup\zeta_j),0)$ is a Gram entry about the coordinate plane for the sources $\gamma^{-1}\cdot\zeta_i$, supported in the corresponding coordinate half-space. (iv) is \cref{prop:s08:OS3} for $\ZKinf(\cdot,0)$ (through \cref{lem:s08:OS2K}; charge conjugation gives the vanishing for odd $n$). (v) is \cref{thm:s09:S2K}(a), including the local convergence of the lattice laws along the full sequence, with \cref{rem:s09:S2K} for the unlocalised fields; (vi) is \cref{thm:s09:S2K}(b) and (c), whose $O(3)$ input is (ii)'s Euclidean clause; the disclaimer on odd moments is that of \cref{thm:s09:S2K}(c).

\emph{C.} Under \Tthree, \cref{prop:s10:U} gives \Itwo\ with amplitude $s_1\in(0,\sigma_1)$ and constant $C_U$ from the torus inputs \cref{hyp:s03:T1,hyp:s03:T2}; for $a\ne z$ the difference $z_k-a_k$ converges by admissibility, so $\sup_k|z_k-a_k|<\infty$ and $C_U$ is replaced as stated. \Ithree\ is \cref{prop:s10:RP}, and the correlation inequalities are \cref{lem:s02:BFS}. So \cref{thm:s11:W} applies: the bound \eqref{eq:s11:W} holds uniformly in $k\ge0$, $J\ge J_1(R)$, with the stated dependence of the constants and no uniformity in $\lmu$. \emph{(i)}: \cref{thm:s11:W}(i) gives the fixed-cube limits $G_J=\lim_kZ^a_{k,J}$ locally uniformly on $\Dtube$ (local boundedness by \eqref{eq:s11:W}, convergence at the real points of the tube by \cref{hyp:s03:W1}, Vitali on the connected tube), and \cref{lem:s11:TWtube}(iii) the limit $\Ginf=\lim_JG_J$, locally uniformly on $\Dtube$, holomorphic, with the bound in the limit; \cref{thm:s11:W}(ii) gives the Schwinger bounds, the factor $2^{n+p}$ and \Ezz. On the Wick tube $T_{\cF'}$ of \cref{def:s08:enlarged}, \cref{prop:s08:T}(a) gives $Z^a_{\infty,T}$, equal to $\Ginf$ on $T_{\cF'}\cap\Dtube$ as the same $J$-limit of the same $G_J$. \emph{(ii)}: translations are \cref{lem:s11:TWtube}(v). For the cubic group, \cref{prop:s12:ID}(c) shows that for an admissible $a$ under \Tthree\ the hypotheses \Bcell\ (\cref{hyp:s03:W1} at the one cell and source of \cref{hyp:s12:Bcell}) and \SCH\ (\cref{hyp:s12:SCH}) hold, every increasing sequence of spacings is a continuum subsequence for the Wick sector (\cref{def:s12:contsub}), and the LIM state $Z^{a,\cK}_\infty$ of every continuum subsequence $\cK$ is $\Ginf$ on $\Dtube$; the Wick sector of \cref{thm:s12:LIM} under \Tthree, \Bcell, \SCH\ gives $Z^{a,\cK}_\infty(\tau_bR\cdot\zeta)=Z^{a,\cK}_\infty(\zeta)$ on the tube with $W^a$ a scalar. Hence $\Ginf$ is $\Cub$-invariant on $\Dtube$, and the $S_{n,p}$ are by differentiation at $0$. \emph{(iii)}: \cref{thm:s08:OS2WT} under \cref{hyp:s03:W1}, \Itwo, \Ithree\ and \cref{lem:s02:BFS}, whose Gram entries exist by \cref{prop:s08:T}(a) and whose passage to the planes $c\notin D$ is \cref{prop:s08:T}(b); the polynomial form is \cref{prop:s08:poly}. \emph{(iv)} is \cref{prop:s08:OS3} for $Z^a_{\infty,T}$. \emph{(v),(vi)}: the Wick sector of \cref{thm:s12:LIM}, clauses (v) and (vi), for the LIM state of the full sequence, which is $\Ginf$ by \cref{prop:s12:ID}(c); the lattice-law clause then holds along the full sequence (\cref{cor:s12:K}). The Gram points of (v)'s characteristic functional have purely imaginary coefficients and lie in the tube, and the reflection positivity consumed there is (iii) at amplitude $0$ (\cref{rem:s13:amplitudes}). For the scheme clause, the shift constant of \cref{thm:s12:LIM} is $\lim_k(a'_k-a_k)$, which exists by admissibility, so no further subsequence is needed and it is independent of $\cK$ (\cref{cor:s12:K}).

\emph{D.} \Cref{thm:s13:TWK,prop:s13:WK}.
\end{proof}

\begin{remark}[the Euclidean clause]\label{rem:s13:rotation}
The Euclidean clauses of B(ii), (iii), (vi) rest on the rotation theorem \cref{hyp:s03:RC}, proved in the rotation part \moving{rotation part}, and on the identification of its limit state with $\ZKinf$ through \cref{cor:s04:G2}; the Wick sector carries no Euclidean clause (\cref{thm:s23:FormA}).
\end{remark}

\begin{theorem}[the Wick tube on King's convention]\label{thm:s13:TWK}
Let the couplings satisfy \Tthree\ and let $a$ be an admissible scheme. Let $\sigma_+=s_1$, $H$, $\kappa$ be the constants of \eqref{eq:s11:constants} and let $K_-$ be the constant there, computed from the cell constant $B^{(a)}_{\rm cell}$ of the comparison convention (\cref{hyp:s12:Bcell}), which is finite. Then for every finite family of real sources with supports in $Q_R$, every $k\ge0$, every $J\ge J_1(R)$ and every $\zeta\in\Dtube$,
\begin{equation}\label{eq:s13:WK}
|Z^{a,K}_{k,J}(\zeta)|\ \le\ Z^{a,K}_{k,J}(\Re\zeta)\ \le\ \exp\big(\norm{g(\zeta)}_{2,k}^2+\kappa\norm{k_+(\zeta)}_{1,k}+K_-\wnorm{k_-(\zeta)}\big),
\end{equation}
the bound \eqref{eq:s11:W} for the K-family functionals with the constants of the comparison convention; and the fixed-cube limits $G^K_J(\zeta)=\lim_{k\to\infty}Z^{a,K}_{k,J}(\zeta)$ exist for every $\zeta\in\Dtube$, locally uniformly, so that \cref{thm:s11:W}(i),(ii) hold on $\NDK$: $G^K_J$ and $G^K_\infty=\lim_JG^K_J$ exist locally uniformly on $\Dtube$ and are holomorphic, the bound \eqref{eq:s13:WK} passes to the limits, and the Schwinger bounds and \Ezz\ of \cref{thm:s13:main}C(i) hold for $G^K_\infty$.
\end{theorem}

\begin{proof}[Sketch]
The proof of \cref{thm:s11:W} is run on $\NDK$; every step of it that sees the states only as members of the class of \cref{lem:s02:BFS} is unchanged, and three steps are convention-sensitive.

(1) \emph{The negative leg and its direction.} By \cref{lem:s04:E} every source-carrying cell of the wall insertion of \cref{lem:s11:W3} has all its sites alive, so \cref{lem:s11:W4,lem:s11:W5,lem:s11:W67,prop:s11:W8} hold for $\avK{\cdot}$ as printed, with $K_-$ computed from the cell constant of the comparison convention (\cref{rem:s11:Kconvention}). That constant is finite: killing the sites of a cell increases a class-$\cD$ functional (\cref{lem:s04:E}(e)), so $B^{(a)}_{\rm cell}\le B^K_{\rm cell}$, and $B^K_{\rm cell}<\infty$ by \cref{hyp:s03:W1} at the one cell and the one negative source of \cref{def:s11:cells}.

(2) \emph{The positive leg.} The interlacing inequality $\avK{\cdot}\le\ava{\cdot}$ on class $\cI$ (Step~1 of \cref{lem:s04:G}, \cref{rem:s04:F}), prepended to \cref{prop:s11:W1prime}, gives $\ln N^K_{k,J}(0,q)\le\kappa\norm q_{1,k}$ for every lattice function $0\le q\le s_1$ vanishing on $\TJ$ (\cref{cor:s11:W2prime} on $\NDK$); \Itwo\ comes from \Tthree\ and the torus inputs through \cref{prop:s10:U}, \Ithree\ is \cref{prop:s10:RP} and \Ifour\ is \cref{lem:s02:BFS}, none of which sees the cube convention.

(3) \emph{Assembly and limits.} The splitting of the Wick source into its positive and negative parts and the H\"older step of the proof of \cref{thm:s11:W} use only the correlation inequalities and run unchanged, giving \eqref{eq:s13:WK} at every $(k,J)$ with the sources' supports at distance $>1$ from $\R^3\setminus\CJ$, so that the positive part vanishes on $\TJ$. The family $(Z^{a,K}_{k,J})_k$ is then locally bounded on the connected tube and converges at the real points of the tube by \cref{hyp:s03:W1}, hence locally uniformly (Vitali), as in \cref{thm:s11:W}(i); the $J$-limit, the holomorphy and the Schwinger bounds follow as in \cref{prop:s13:WK}(a) below. Kill steps outside the scope of \cref{lem:s04:G}, for an observable that does not depend on the killed sites, are the monotonicity of \cref{lem:s02:BFS} under the replacement of a single-site measure by a point mass.
\end{proof}

\begin{proposition}[the Wick sector on King's convention]\label{prop:s13:WK}
Let the couplings satisfy \Tthree\ and let $a$ be an admissible scheme. Then, with the constants of \cref{thm:s13:main}C:
\begin{enumerate}[label=\textup{(\alph*)},leftmargin=*,nosep]
\item $G^K_J\defeq\lim_kZ^{a,K}_{k,J}$ exists locally uniformly on $\Dtube$ for every $J\ge J_1(R)$, $G^K_\infty\defeq\lim_JG^K_J$ exists locally uniformly on $\Dtube$, both are holomorphic, $|G^K_\infty(\zeta)|\le\exp(\norm{g(\zeta)}_2^2+\kappa\norm{k_+(\zeta)}_{L^1}+K_-\wnorm{k_-(\zeta)})$, and the Schwinger bounds and \Ezz\ of \cref{thm:s13:main}C(i) follow.
\item $G^K_\infty(\tau_bR\cdot\zeta)=G^K_\infty(\zeta)$ on $\Dtube$ for all $b\in\R^3$, $R\in\Oh$, $W^a$ a scalar.
\item Reflection positivity about every coordinate plane in the matrix form at amplitude $<\sigma_W$ and in the polynomial form, and \OSth, for $G^K_\infty$ and its Schwinger functions.
\item The measure and weak \OSf\ of \cref{thm:s13:main}C(v),(vi) for the joint field under $\NDK$, with the local convergence of the laws of $\chi(F_{\Lk},W^a_{\Lk})$ under $\NDK_{k,J}$ along the full sequence $k\to\infty$ at every fixed $J\ge\max(J_0(\chi),J_1(R_\chi))$, then $J\to\infty$.
\end{enumerate}
\end{proposition}

\begin{proof}
The proofs of \cref{thm:s11:W}(i), \cref{lem:s11:TWtube}, \cref{thm:s06:S1K} with scalar sources, \cref{thm:s08:OS2WT} with \cref{prop:s08:T}, and \cref{thm:s09:S2K} in the Wick sector are run on $\NDK$ with the replacements below; every other step of theirs is convention-free, since it sees the states only as members of the class of \cref{lem:s02:BFS} or concerns the limit object alone.

(a) By \cref{thm:s13:TWK} the bound \eqref{eq:s11:W} holds for $Z^{a,K}_{k,J}$ uniformly in $k\ge0$, $J\ge J_1(R)$, and $\lim_kZ^{a,K}_{k,J}(\zeta)$ exists for $\zeta\in\Dtube$; so $(Z^{a,K}_{k,J})_k$ is locally bounded and convergent on the connected tube, and $G^K_J$ exists locally uniformly and is holomorphic (Vitali), with the bound. The $J$-limit is the argument of \cref{lem:s11:TWtube}(iii), whose one convention-sensitive input is the monotonicity of the cone functionals in the cube; on $\NDK$ that is \cref{cor:s04:G1} (for $J\ge J_1(R+\frac12)$, which exceeds the $J_0$ of that corollary). The split family, the uniqueness set, Vitali on the split tube and the extension to $\Dtube$ are as there; the Schwinger bounds are Cauchy's estimates on a polydisc in the tube, as in \cref{thm:s11:W}(ii).

(b) The steps of \cref{thm:s06:S1K} with scalar sources. Coordinate permutations are exact on $\NDK$ for $(f,h)$ by \cref{lem:s06:PK}. The reflection sandwich is \cref{lem:s04:C}(b) at the cone points of the tube: for $(f,g)\in\cK$ with $\norm g_\infty<\sigma_+/2$ the point $(\mathrm i,1)$ lies in $\Dtube$, so (a) supplies the three $k$-limits and the two $J$-limits of the two-step sandwich $C^K_{J\mp2}$. Conjugation and group closure on the cone are unchanged. Off the cone, the identity theorem on the split family's tube of \cref{lem:s11:TWtube}, where (a) gives the limits, extends the identity from the cone points to the connected tube; so \cref{lem:s04:C} is needed on $\cK$ only, and no class-$\cD$ form of it enters (the proof splits $h=h^+-h^-$ and never applies the sandwich to a negative scalar source). Translations are \cref{lem:s06:TK}, stated on $\cK$ and applied at the cone points of the tube with (a) for the limits, then the continuity route of \cref{cor:s06:TK} with the Wick-sector second-moment bound, which Cauchy's estimates give from \eqref{eq:s11:W} on $\NDK$; the translated family needs $J\ge\max(J_0+1,J_1(R+M+\frac12))$. Invariance about every point and differentiation concern the limit object.

(c) \Cref{thm:s08:OS2WT} and \cref{prop:s08:T} on $\NDK$: in the sandwich of the $J$-limit the upper box is \cref{lem:s04:D} (the printed boxes $\CJ$, $C_{J+2n}$) and the lower end carries the kill step of \cref{rem:s04:F}; the positive-leg bound \cref{cor:s11:W2prime} holds on $\NDK$ by \cref{rem:s04:F} prepended to \cref{prop:s11:W1prime} (\cref{rem:s11:Kconvention}); the cube monotonicity in both classes is \cref{cor:s04:G1}; the limits are (a) and the translations to the planes $c\notin D$ are (b). \OSth\ and the polynomial form use only the existence and holomorphy of the limit near $0$ and charge conjugation (\cref{prop:s08:OS3,prop:s08:poly}).

(d) The steps of \cref{thm:s09:S2K} on $\cV=\Lambda^2\oplus\R$ with the scheme-shifted field: \eqref{eq:s11:W} on $\NDK$ in the moment bounds, (a) for the limit functionals and the lattice laws, (b) and (c) for the invariance and the reflection positivity consumed by the abstract chain of \cite{SymPhase}. The second moments of the lattice field under $\NDK_{k,J}$ come from \eqref{eq:s11:W} on $\NDK$ by Cauchy's estimates, so tightness and the identification of the limit law run as in the curvature sector, with the fixed-cube limits of (a) along the full sequence, for $J\ge\max(J_0(\chi),J_1(R_\chi))$; the remaining steps concern the limit measure.
\end{proof}

\begin{remark}[the three amplitude constants]\label{rem:s13:amplitudes}
Three amplitudes occur in the Wick sector, all built from $s_1$: the matrix form of \OSt\ is proved at amplitude $\amp(\zeta_i)<\sigma_W=\sigma_+/8$ on the Wick tube (\cref{thm:s08:OS2WT}); the tube $\Dtube$ has half-widths $\sigma_+/2$ for the positive and $H/2=\sigma_+/32$ for the negative part of the real Wick source, and the form of \OSt\ proved for LIM states (\cref{thm:s12:LIM}) takes $\amp(\zeta_i)<H/2<\sigma_W$, so that every Gram point lies in $\Dtube$, where those limits exist; and $\sigma_0=H/4$ is the radius of the polydisc in Cauchy's estimates for \Ezz. The consumer of \OSt\ in clauses (v),(vi), the reflection positivity of exponentials $e^{\mathrm i\Phi(u)}$ with real $u$ in the clustering theorems of \cite{SymPhase}, needs purely imaginary Wick coefficients only, that is amplitude $0$ (\cref{rem:s08:amplitudes}).
\end{remark}

\begin{remark}[the photon mass and the couplings]\label{rem:s13:mu0}
The photon mass $\lmu>0$ is fixed throughout \cref{thm:s13:main}. No relation $\lmu\le m$ is used or needed: the one input whose printed proof in the torus paper \cite{SymPhase} carries the restriction $\lmu\in(0,\mus]$, the lattice reflection positivity of King's tori, holds for all $\lmu>0$, since that proof uses no bound on $\lmu$ and the paper's own lemma on the finite-volume photon covariance lifts the scope (\cref{prop:s10:RP}, \cref{lem:s14:CW}(d)). Reflection positivity of the cube limit needs the margin $R(\lmu)=\pi\sqrt3/(2\lmu)$ of \eqref{eq:s07:margin} inside the upper half of the sandwich (\cref{lem:s07:U}) and nowhere in the conclusion. The constants of Leg~C are absolute; those of Leg~W depend on $\lmu$ and are not uniform in it (a tube bound uniform in the photon mass is the first step of the separate work on the limit of vanishing photon mass). The couplings range over all of $e>0$, $\lambda,m,\lmu>0$: their quantifier is the sources' (\cref{lem:s14:CW}(e), \cref{rem:s14:quantifier}), the stopping scale of the lattice constructions behind \cref{hyp:s03:EC,hyp:s03:W1,hyp:s03:T1} being chosen after the point and the datum, with the mass cap of \cite[Prop.~3.1]{Ba2} one more one-sided cap on that scale (\cref{rem:s14:BALm}); so the theorem contains no sentence ``$e,\lambda$ sufficiently small'', and the only coupling condition of the Wick sector is the one-sided sign clause \Tthree\ (\cref{lem:s14:CW}(b)), proved on the explicit regions of \cref{prop:s14:T3}.
\end{remark}
